\documentclass[11pt]{article}

\usepackage[letterpaper,margin=0.82in]{geometry}
\usepackage[T1]{fontenc}
\usepackage{amsmath}
\usepackage{booktabs}
\usepackage{float}
\usepackage{graphicx}
\usepackage{microtype}
\usepackage[numbers,sort&compress]{natbib}
\usepackage[capitalize]{cleveref}
\usepackage{xcolor}
\usepackage{enumitem}
\usepackage{tabularx}

\setlength{\emergencystretch}{3em}
\setlength{\bibsep}{1pt}
\setlength{\textfloatsep}{10pt plus 2pt minus 2pt}
\setlength{\floatsep}{9pt plus 2pt minus 2pt}
\setlength{\intextsep}{9pt plus 2pt minus 2pt}
\renewcommand{\topfraction}{0.95}
\renewcommand{\textfraction}{0.05}
\renewcommand{\floatpagefraction}{0.75}
\pagestyle{plain}

\definecolor{statusred}{HTML}{9F1239}
\definecolor{statusfill}{HTML}{FFF1F2}
\newcommand{\draftwarning}[1]{%
  \begin{center}
  \fcolorbox{statusred}{statusfill}{%
    \begin{minipage}{0.93\linewidth}\small\textbf{Study status.} #1\end{minipage}}
  \end{center}}

% Generated by code/build_pilot_artifacts.py; do not edit by hand.
\newcommand{\PilotEvidenceStatus}{\texttt{clean\_v2\_retrospective\_estimation}}
\newcommand{\PilotNModels}{4}
\newcommand{\PilotNSeeds}{5}
\newcommand{\PilotNRepresentedBenchmarks}{3}
\newcommand{\PilotNTransferBenchmarks}{4}
\newcommand{\PilotNAPDatasets}{7}
\newcommand{\PilotNTransferPositiveVsBase}{2}
\newcommand{\PilotNSeedPairs}{10}
\newcommand{\PilotBaseRepresented}{0.658}
\newcommand{\PilotBaseTransfer}{0.866}
\newcommand{\PilotSFTRepresented}{0.982}
\newcommand{\PilotSFTTransfer}{0.807}
\newcommand{\PilotCompositionRepresented}{0.962}
\newcommand{\PilotCompositionTransfer}{0.883}
\newcommand{\PilotLogitTransfer}{0.891}
\newcommand{\PilotRepDeltaSFT}{-0.019}
\newcommand{\PilotRepDeltaSFTBootstrapMean}{-0.019}
\newcommand{\PilotRepDeltaSFTLow}{-0.031}
\newcommand{\PilotRepDeltaSFTHigh}{-0.010}
\newcommand{\PilotTransferDeltaSFT}{+0.076}
\newcommand{\PilotTransferDeltaSFTBootstrapMean}{+0.075}
\newcommand{\PilotTransferDeltaSFTLow}{+0.058}
\newcommand{\PilotTransferDeltaSFTHigh}{+0.093}
\newcommand{\PilotTransferDeltaBase}{+0.017}
\newcommand{\PilotTransferDeltaBaseBootstrapMean}{+0.017}
\newcommand{\PilotTransferDeltaBaseLow}{+0.005}
\newcommand{\PilotTransferDeltaBaseHigh}{+0.030}
\newcommand{\PilotTargetFPRPct}{5}
\newcommand{\PilotBaseTransferFPRPct}{8.1}
\newcommand{\PilotSFTTransferFPRPct}{15.5}
\newcommand{\PilotCompositionTransferFPRPct}{11.4}
\newcommand{\PilotConvexWeight}{0.95}
\newcommand{\PilotLockHashShort}{cabc8dee}
\newcommand{\PilotScoreHashShort}{b941ddba}
\newcommand{\PilotCompositionHashShort}{92c2cbc3}


\title{Compose, Don't Tune?\\
\large Recovering Held-Out Transfer by Averaging Base and Fine-Tuned Guard Scores --- A Retrospective Fixed-Panel Pilot}
\author{Reza Rahimi, PhD\\
  \small JazzX AI, Los Altos, CA, USA\\
  \small\texttt{reza.rahimi@jazzx.ai}}
\date{Working draft --- July 2026}

\begin{document}
\maketitle

\draftwarning{This manuscript is a research draft built from \emph{retrospective} clean-v2
evidence.  The transfer benchmarks were withheld from SFT but inspected during method
development.  No separately locked prospective Paper B study has run.  Sections describing
future validation are protocol, not results.}

\begin{abstract}
Fine-tuning a compact language model into a binary safety guard can improve discrimination on
benchmark sources represented during adaptation while weakening transfer to other datasets.  We
ask whether keeping the same checkpoint's pre-adaptation decision in an output-space composition (that is, by averaging the two models' final calibrated scores rather than merging their internal weights)
can recover transfer relative to its supervised fine-tuning (SFT) adapter.  Using the lock-verified per-row scores from a
companion study, we retrospectively compare \PilotNModels{} unadapted instruction checkpoints,
\PilotNSeeds{} LoRA-SFT adapters per checkpoint, and a fixed average of separately calibrated
unsafe probabilities.  Relative to SFT, the observed composition contrast is
\PilotRepDeltaSFT{} in represented-source macro average precision (AP: ranking quality measured as the area under the precision--recall curve, with ``macro'' meaning averaged equally across benchmarks), with a descriptive
paired-bootstrap 95\% interval [\PilotRepDeltaSFTLow, \PilotRepDeltaSFTHigh], and
\PilotTransferDeltaSFT{} on dataset-held-out transfer
([\PilotTransferDeltaSFTLow, \PilotTransferDeltaSFTHigh]).  The transfer contrast relative to the
base is heterogeneous: positive for \PilotNTransferPositiveVsBase{} checkpoints, near zero for one, and negative for
Qwen3-4B.  Thresholds chosen for a \PilotTargetFPRPct\% calibration false-positive-rate (FPR)
target yield \PilotCompositionTransferFPRPct\% transfer macro-FPR for the composition, so rank
recovery does not imply calibration transfer.  The method also requires two inference passes.
Because the operator and transfer results informed this analysis, these are conditional
fixed-panel retrospective estimates, not confirmatory evidence.  They motivate a separately
locked prospective comparison against WiSE-FT, an equal-inference-cost SFT+SFT ensemble, and
matched-compute preservation tuning.
\end{abstract}

\begin{center}
\fbox{\begin{minipage}{0.95\linewidth}\small
\textbf{Reader's guide (plain-language terms used throughout).}
\begin{description}[leftmargin=1.2em,style=nextline,itemsep=1pt,topsep=2pt]
  \item[Base / SFT adapter.] The \emph{base} is the instruction-tuned checkpoint before we train it into a guard; the \emph{SFT adapter} is that same checkpoint after supervised fine-tuning (SFT) on labeled safe/unsafe prompts, done with LoRA.
  \item[Unsafe score / calibration.] Each guard outputs a number meaning ``how unsafe.'' \emph{Calibration} rescales that number on a held-out development split so it reads as a probability that is comparable across guards.
  \item[Output-space composition.] Combining two guards by averaging their (calibrated) output scores, rather than by averaging their internal weights; it needs only comparable scores but runs both models.
  \item[Represented-source vs.\ dataset-held-out (transfer).] \emph{Represented} benchmarks contributed rows to SFT; \emph{transfer} benchmarks did not, so they test generalization to other datasets.
  \item[AP / macro-AP.] Average precision (AP) summarizes ranking quality across all cutoffs (area under the precision--recall curve); \emph{macro-AP} averages AP equally over benchmarks.
  \item[Operating point / FPR / TPR.] Picking one cutoff turns scores into decisions; FPR is the fraction of safe prompts wrongly flagged, TPR the fraction of unsafe prompts caught.
  \item[Estimand.] The exact quantity we set out to estimate (here, a difference in macro-AP), stated before looking at the result.
  \item[Paired hierarchical bootstrap / percentile interval.] Re-sampling the data many times---keeping the two guards paired on the same rows and respecting the checkpoint$\to$seed grouping---to get a 95\% range for a contrast.
  \item[Noninferiority margin ($m$).] A small, pre-chosen allowed drop: ``no worse than the reference by more than $m$.''
  \item[LCB.] Lower confidence bound---the low end of a one-sided interval.
  \item[clean-v2 / lock / hash.] ``clean-v2'' is the reviewed Paper~A data snapshot used here; a \emph{lock}/\emph{hash} is a fingerprint of the exact inputs, so every table can be traced back to specific data.
\end{description}
\end{minipage}}
\end{center}

\section{Introduction}
\label{sec:introduction}

Parameter-efficient adaptation makes it inexpensive to specialize an instruction model into a
prompt-safety guard.  LoRA, for example, updates a low-rank parameterization rather than all model
weights \citep{hu2021lora}.  Yet specialization creates a practical tension: a guard can become
better at the policies and benchmark sources represented during SFT while losing useful ranking
information on dataset-held-out sources. Here \emph{represented sources} are the datasets whose rows were used to train the guard, while \emph{dataset-held-out} (or \emph{transfer}) sources are other datasets whose rows were never used for training, so they test how well the guard generalizes.  The companion Paper A study measures this tension under
a controlled base-to-adapter design \citep{rahimi2026benchmark}.  This draft asks the next, narrower
question: can a fixed composition recover some transfer relative to SFT without another training
run?

The candidate rule is deliberately simple.  For one input, run the unadapted checkpoint and one
of its SFT adapters, calibrate their unsafe scores using development calibration rows, and average
the two calibrated probabilities with fixed equal weights.  Here \emph{base} means unadapted to
the guard SFT recipe; it does not mean pretrained from scratch, unaligned, or generally untuned.
The rule is output-space and therefore does not require the two systems to expose interpolable
weights, but it does require aligned labels, compatible unsafe-score semantics, separate
calibration, and two inference passes.

The current evidence is a feasibility result, not the final Paper B study.  Paper A's score table
and lock establish what inputs were analyzed, but they do not prospectively bind a composition
operator invented later.  Moreover, some transfer benchmarks and earlier results were seen during
development.  We therefore use the current panel to estimate a retrospective retention--recovery
trade-off, expose failure modes, and write a protocol that can later be frozen before a genuinely
uninspected cohort is evaluated.

This draft makes three bounded contributions:

\begin{enumerate}[leftmargin=1.5em]
  \item a generated, provenance-gated summary of the clean-v2 retrospective composition pilot
        (every reported number is auto-filled from a locked results file, so the prose cannot
        silently drift from the underlying data);
  \item an explicit claim boundary separating score composition from prior weight-space and
        calibrated-ensemble methods (a clear statement of what is and is not new relative to
        averaging model \emph{weights} or calibrating existing ensembles); and
  \item an executable prospective-protocol contract that fails claim-bearing validation until
        the cohort, margin, controls, statistics, systems measurements, and software lock are fixed
        (a machine-checked checklist that refuses to certify the study as confirmatory until every
        required piece is in place).
\end{enumerate}

It does \emph{not} claim Pareto dominance (being at least as good on \emph{every} axis and strictly better on at least one), universal improvement, a causal mechanism, a validated
base-competence rule, or deployment readiness.

\begin{center}
\fbox{\begin{minipage}{0.93\linewidth}\small
\textbf{What is new here (and what is not).}
\emph{The invention} is one fixed, no-retraining decision rule for safety guards: on the same input,
run the model \emph{before} guard fine-tuning (the \emph{base}) and one guard-fine-tuned \emph{adapter},
convert each model's raw unsafe score into a probability using a held-out development split
(\emph{calibration}), and report the plain equal-weight average of the two probabilities---fixed at
$\tfrac12$ and $\tfrac12$---as the guard's score (\Cref{eq:composition}). Nothing is retrained and no
weight is chosen from the test rows.
\emph{Why it might help:} guard fine-tuning tends to sharpen ranking on the data sources it was trained
on while eroding ranking on unseen sources; keeping the base inside the average retains some of that
unseen-source signal.
\emph{What is genuinely new} is \emph{not} score averaging itself---calibrated ensembles, weight-space
interpolation, and weight averaging already exist
\citep{kumar2022calibrated,wortsman2022wiseft,wortsman2022modelsoups}---but its use as a
retraining-free rule for recovering the base-versus-fine-tuned safety-guard trade-off documented in the
companion study \citep{rahimi2026benchmark}, together with an explicit claim boundary and a
machine-checked prospective-validation contract.
\emph{What is not claimed:} not that the average beats both of its parts everywhere, not a proven
mechanism, not a validated when-to-use rule, and not deployment readiness. The evidence in this draft is
retrospective and estimation-only.
\end{minipage}}
\end{center}

\section{Related work and novelty boundary}
\label{sec:related}

The single most closely related work is the companion Paper~A study, which under a controlled base-to-adapter design documents the specialization trade-off this draft tries to counteract: guard SFT can sharpen represented-source ranking while eroding dataset-held-out transfer \citep{rahimi2026benchmark}. That study supplies the fixed panel and the row-keyed base and SFT scores reused here, but it does not itself propose or evaluate any score-composition operator, so it establishes the problem rather than the remedy examined below.

The underlying ensemble idea is not new.  Calibrated ensembling has been shown to mitigate
in-distribution and distribution-shift trade-offs in image classification when component models
are calibrated using in-distribution data \citep{kumar2022calibrated}.  WiSE-FT interpolates the
weights of a zero-shot model and its fine-tuned descendant to improve robustness under shift
\citep{wortsman2022wiseft}.  Model soups average compatible fine-tuned weights and can avoid the
runtime cost of a conventional multi-pass ensemble \citep{wortsman2022modelsoups}.

In plain terms, these prior lines all \emph{combine} models, but in different ways. Calibrated ensembles average the predictions of separately and \emph{independently} trained models \citep{kumar2022calibrated}; WiSE-FT and model soups combine models in \emph{weight space}, blending internal parameters into a single network that then runs in one pass \citep{wortsman2022wiseft,wortsman2022modelsoups}. WiSE-FT is the closest analog to the rule studied here, because it likewise combines a model with its own fine-tuned descendant---the same base-and-its-own-adapter pairing used below---but it averages \emph{weights}, whereas this paper averages the two models' final calibrated decision scores in \emph{output space}. Output-space blending is more portable (it works even when the two networks' weights cannot be averaged, e.g.\ across different architectures) but pays for that flexibility by running both models on every input.

This study does not present score averaging as a novel primitive.  Its narrower target is a safety-
guard setting in which a single guard SFT may improve represented-source ranking while degrading
dataset-held-out ranking.  Output-space composition is architecture-agnostic in the limited sense
that it needs comparable decision scores rather than interpolable parameters.  That portability
comes with a systems cost: both components must run.  A submission-grade study must therefore
compare against actual WiSE-FT rescoring, not merely discuss it, and must measure whether an
SFT+SFT two-pass ensemble obtains the same benefit.  Without the latter control, any gain could be
generic variance reduction from a second model pass rather than information uniquely retained by
the base.

The concern in plain terms: averaging \emph{any} two models tends to smooth out noise and help a little. So the base+SFT average must be checked against an average of two \emph{independently fine-tuned} adapters (the SFT+SFT control). If SFT+SFT helps just as much, the gain is generic ``two-model'' averaging rather than anything special about keeping the base.

\section{Evidence status and research questions}
\label{sec:status}

\Cref{tab:evidence-tiers} separates the evidence tiers used in this draft.  The distinction is
substantive: clean execution provenance does not erase prior exposure to evaluation sources, and a
Paper A lock cannot bind Paper B code written after it.

\begin{table}[H]
\centering
\caption{Evidence tiers.  Only the first two rows exist; the prospective row is planned.}
\label{tab:evidence-tiers}
\small
\begin{tabularx}{\linewidth}{@{}lXl@{}}
\toprule
Tier & What it establishes & Current status \\
\midrule
Paper A clean-v2 scores & Row-keyed base and SFT scores bound to the Paper A execution contract & Complete \\
Retrospective composition & Reproducible post-lock analysis of fixed, previously exposed transfer benchmarks & \PilotEvidenceStatus \\
Prospective Paper B & Frozen operator, margins, controls, model panel, uninspected cohort, and analysis & Not executed \\
\bottomrule
\end{tabularx}
\end{table}

The pilot addresses three research questions, all conditional on the fixed \PilotNModels{}-checkpoint panel:

\begin{description}[leftmargin=3.5em,style=nextline]
  \item[RQ1: represented retention.] What is the composition-minus-SFT change on sources
  represented during SFT?  The pilot estimates the loss; it does not label that loss
  noninferior because no margin was prospectively justified.

  \item[RQ2: transfer recovery.] What is the composition-minus-SFT change on dataset-held-out
  transfer benchmarks, and how does the composition compare descriptively with the base?

  \item[RQ3: heterogeneity hypothesis.] Can base competence measured on a locked development
  cohort predict composition-minus-SFT recovery on a disjoint prospective cohort?  The current
  \PilotNModels{} points motivate this question but cannot validate a decision rule.
\end{description}

\section{Retrospective pilot methods}
\label{sec:methods}

\subsection{Panel, data roles, and score artifact}

The pilot reuses Paper A's fixed panel: Qwen2.5-1.5B, SmolLM2-1.7B, SmolLM3-3B, and Qwen3-4B,
with \PilotNSeeds{} LoRA-SFT seeds per checkpoint.  Evaluation sources are grouped into represented-source
tests and dataset-held-out transfer tests.  Dataset-held-out means those rows were not used for
SFT; it does not mean the benchmark identities or earlier outcomes were unseen by the researcher.
All absolute and contrast values consumed by this manuscript are generated from the canonical
composition JSON rather than copied into the TeX source.

The publication adapter refuses legacy evidence, a nonfinal run, a changed operator, an altered
seed panel, or a prospective flag.  It binds the exact retrospective result with composition hash
\texttt{\PilotCompositionHashShort\ldots}, Paper A lock
\texttt{\PilotLockHashShort\ldots}, and score artifact
\texttt{\PilotScoreHashShort\ldots}.  These hashes make the draft stable; they do not constitute a
prospective Paper B lock.

\subsection{Fixed primary operator}

Let $s_b(x)$ and $s_{a,j}(x)$ denote the unsafe scores for base $b$ and SFT seed $j$.  Let
$C_b$ and $C_{a,j}$ be calibrators fit only on the development calibration split. A \emph{calibrator} is a simple order-preserving rescaling, learned only on a separate development split, that maps a guard's raw score onto a $0$--$1$ probability, so that ``$0.7$'' means roughly the same thing for both systems and the two can be averaged meaningfully.  The fixed
primary composition is
\begin{equation}
  s_{\mathrm{comp},j}(x)
  = \tfrac{1}{2} C_b(s_b(x)) + \tfrac{1}{2} C_{a,j}(s_{a,j}(x)).
  \label{eq:composition}
\end{equation}
In words, \Cref{eq:composition} says: pass each system's raw ``how unsafe'' score through its own calibrator $C$ to get a comparable unsafe probability, then report the plain midpoint of the two. Scoring uses \emph{both} the pre-fine-tuning \emph{base} checkpoint and one \emph{fine-tuned adapter}; nothing is retrained, and the equal $\tfrac{1}{2}$ weights are fixed in advance, so the operator has no free parameter that could be tuned to the reported test rows.
The analyzer fixes equal weights without fitting reported test rows.  Because the rule was formalized
after earlier transfer exposure, this implementation fact is not human test blindness.  The remaining
score combiners are exposed ablations and cannot become primary after observing transfer.

\subsection{Estimands and uncertainty}

For regime $r\in\{\mathrm{represented},\mathrm{transfer}\}$, the primary descriptive estimand is
\begin{equation}
  \Delta_r = \operatorname{macroAP}_r(\mathrm{composition})
             - \operatorname{macroAP}_r(\mathrm{SFT}).
\end{equation}
AP is computed per benchmark with tie-aware ranking, macro-averaged over benchmarks, averaged
over SFT seeds, and then averaged equally across the \PilotNModels{} fixed checkpoints.  The composition-
minus-base contrast is a secondary anchor.  It is not the main success criterion because the goal
is to recover transfer relative to SFT while retaining a bounded amount of represented-source
performance.

Descriptive intervals use a paired family-by-seed hierarchical bootstrap.  Family resampling is
shared across both sides of a contrast; SFT seeds are resampled within checkpoint.  The base has
one scoring pass, so seed resampling captures adapter-side variation but not uncertainty over base
training seeds.  The reported point beside each percentile interval is the observed contrast.  A
separate generated macro retains the bootstrap mean so the two quantities cannot be silently
mixed.

The \PilotNModels{} checkpoints and \PilotNAPDatasets{} AP datasets are treated as a fixed panel, not random samples from
model or benchmark populations.  Intervals summarize row-family and adapter-seed variation
conditional on this panel and protocol; they do not license universal model-family claims or
formal population-level significance language.

\subsection{Operating-point diagnostic}

Macro-AP is threshold-free.  A threshold is selected from development calibration negatives for a
target FPR of \PilotTargetFPRPct\%.  The artifact retains both regimes; \Cref{tab:pilot-operating}
reports transfer macro and pooled FPR/TPR.  This secondary diagnostic is not a production guarantee.

\section{Retrospective results}
\label{sec:results}

% Generated by code/build_pilot_artifacts.py; do not edit by hand.
\begin{table}[H]
\centering
\caption{Retrospective clean-v2 fixed-panel macro-AP. The calibrated average is the fixed primary operator. Logit averaging is an exposed ablation and cannot be promoted after observing transfer results.}
\label{tab:pilot-summary}
\small
\begin{tabular}{lrrr}
\toprule
Guard & Represented & Transfer & $\min$(both) \\
\midrule
Unadapted base & 0.658 & 0.866 & 0.658 \\
SFT adapter & 0.982 & 0.807 & 0.807 \\
Base+SFT calibrated average & 0.962 & 0.883 & 0.883 \\
Base+SFT logit average (ablation) & 0.943 & 0.891 & 0.891 \\
\bottomrule
\end{tabular}
\end{table}


The calibrated composition changes represented-source macro-AP from
\PilotSFTRepresented{} for SFT to \PilotCompositionRepresented{}, an observed difference of
\PilotRepDeltaSFT{} with a paired-bootstrap percentile interval
[\PilotRepDeltaSFTLow, \PilotRepDeltaSFTHigh].  On transfer it changes macro-AP from
\PilotSFTTransfer{} to \PilotCompositionTransfer{}, an observed difference of
\PilotTransferDeltaSFT{} [\PilotTransferDeltaSFTLow, \PilotTransferDeltaSFTHigh].  These estimates
describe a measured exchange: lower represented-source AP than SFT and higher transfer AP than
SFT.  Whether the represented loss is acceptable is a decision question requiring a prospectively
supported noninferiority margin.

The panel-average composition-minus-base transfer contrast is \PilotTransferDeltaBase{}
[\PilotTransferDeltaBaseLow, \PilotTransferDeltaBaseHigh], but the panel mean hides important
heterogeneity.  \Cref{tab:pilot-models} shows \PilotNTransferPositiveVsBase{} positive observed composition-minus-base contrasts
for SmolLM2 and Qwen2.5, an interval spanning zero for SmolLM3, and a negative contrast for
Qwen3-4B.  Thus the pilot supports transfer recovery relative to SFT on these checkpoints; it does
not support a claim that composition beats both components per checkpoint.

% Generated by code/build_pilot_artifacts.py; do not edit by hand.
\begin{table}[H]
\centering
\caption{Transfer macro-AP by checkpoint. Deltas are observed contrasts with descriptive paired-bootstrap percentile intervals. Recovery relative to SFT is positive for every checkpoint, but comparison with base is heterogeneous.}
\label{tab:pilot-models}
\small
\resizebox{\linewidth}{!}{%
\begin{tabular}{lrrrrr}
\toprule
Checkpoint & Base & SFT & Base+SFT & $\Delta$ vs. SFT [95\% CI] & $\Delta$ vs. base [95\% CI] \\
\midrule
SmolLM2-1.7B & 0.790 & 0.830 & 0.857 & +0.027 [+0.013, +0.043] & +0.067 [+0.038, +0.095] \\
Qwen2.5-1.5B & 0.819 & 0.780 & 0.855 & +0.075 [+0.047, +0.103] & +0.036 [+0.008, +0.066] \\
SmolLM3-3B & 0.910 & 0.823 & 0.907 & +0.084 [+0.063, +0.104] & -0.003 [-0.012, +0.005] \\
Qwen3-4B & 0.944 & 0.794 & 0.914 & +0.120 [+0.082, +0.160] & -0.030 [-0.043, -0.018] \\
\bottomrule
\end{tabular}%
}
\end{table}


Logit averaging reaches transfer macro-AP \PilotLogitTransfer{}, above the fixed primary in this
exposed pilot.  It remains an ablation.  Selecting it after observing transfer would convert a
descriptive comparison into an unacknowledged tuning step.  Likewise, the calibration-only convex
search selects SFT weight \PilotConvexWeight{}; its development selection does not make it the
precommitted primary operator.

% Generated by code/build_pilot_artifacts.py; do not edit by hand.
\begin{table}[H]
\centering
\caption{Realized transfer operating points after choosing thresholds for a 5\% FPR target on calibration negatives. Rank recovery does not imply calibration transfer.}
\label{tab:pilot-operating}
\small
\begin{tabular}{lrrr}
\toprule
Guard & Macro TPR & Macro FPR & Pooled FPR \\
\midrule
Unadapted base & 0.517 & 0.081 & 0.043 \\
SFT adapter & 0.581 & 0.155 & 0.170 \\
Base+SFT calibrated average & 0.639 & 0.114 & 0.091 \\
\bottomrule
\end{tabular}
\end{table}


The operating-point result is less favorable than the ranking result.  Against a
\PilotTargetFPRPct\% calibration target, realized transfer macro-FPR is
\PilotBaseTransferFPRPct\% for the base, \PilotSFTTransferFPRPct\% for SFT, and
\PilotCompositionTransferFPRPct\% for the composition.  Composition improves macro-TPR relative
to both components in this diagnostic, but it still misses the FPR target.  Any deployment claim
based only on AP would omit this calibration-transfer failure. Here \emph{macro-FPR} and \emph{macro-TPR} average each benchmark's false-positive and true-positive rate equally, whereas \emph{pooled FPR} lumps all safe prompts together before computing the rate; the gap between the \PilotTargetFPRPct\% target and the realized FPR means the cutoff calibrated on development data does not carry over to the transfer benchmarks.

\section{What the pilot does not establish}
\label{sec:not-establish}

Several tempting interpretations are not supported.

\begin{itemize}[leftmargin=1.5em]
  \item \textbf{Not Pareto dominance.} Composition is lower than SFT on represented-source AP,
  and it is lower than base on transfer for Qwen3-4B.  A supported noninferiority margin can define
  bounded retention; it cannot retroactively create literal Pareto improvement.

  \item \textbf{Not a base-competence law.} Regressing same-cohort base AP against
  composition-minus-base embeds the predictor in the outcome with a negative sign.  That
  mathematical coupling can manufacture an association.  The prospective protocol instead uses
  development-cohort competence to predict a composition-minus-SFT outcome on a disjoint cohort.

  \item \textbf{Not a mechanism.} Signal-destroying and within-class shuffle diagnostics can
  eliminate some explanations but cannot identify why composition changes ranking.

  \item \textbf{Not evidence that keeping the base is uniquely helpful.} The current pilot lacks
  the independently seeded SFT+SFT two-pass control.  The effect may be generic ensembling.

  \item \textbf{Not free.} Output composition requires two inference passes and associated
  latency, throughput, memory, energy, and serving complexity.  The pilot does not yet contain
  standardized systems measurements.
\end{itemize}

\section{Prospective validation protocol --- not yet executed}
\label{sec:prospective}

The repository contains a machine-validated draft contract, not a completed preregistration.  A
claim-bearing run must use a separate Paper B lock that binds the imported Paper A evidence,
prospective cohort, source tree, analysis configuration, environment, RNGs, operator hierarchy,
model panel, estimands, multiplicity policy, and systems protocol.  The validator's locked mode is
designed to fail until every item below is concrete.

\subsection{Cohort and blindness}

The primary outcome cohort must be selected and content-hashed before inspection of its labels or
model scores.  It must be genuinely disjoint from the cohort used to estimate base competence.
Calling a known public benchmark ``held out'' is insufficient if its earlier scores influenced
operator development.  Retrospective and prospective estimates will be reported separately rather
than pooled.

\subsection{Joint success rule}

Before evaluation, the protocol must justify and freeze a represented-source noninferiority
margin $m$.  A prospective success rule should require both
\begin{align}
  \operatorname{LCB}(\Delta_{\mathrm{represented}}) &> -m, \\
  \operatorname{LCB}(\Delta_{\mathrm{transfer}}) &> 0,
\end{align}
under pre-specified one-sided intervals and a multiplicity policy. Here $\operatorname{LCB}(\cdot)$ is the \emph{lower confidence bound}: the low end of a one-sided interval. In words, the rule requires that---with high confidence---represented-source AP falls by less than the margin $m$, \emph{and} that transfer AP genuinely rises above zero. Demanding \emph{both} at once (an intersection rule) is deliberately stricter than checking either one alone, and the multiplicity policy just fixes how the two tests are combined so the joint error rate stays controlled.  The value of $m$ cannot be
chosen from the observed pilot loss.  Even if both criteria pass, the claim should be bounded
retention plus transfer recovery relative to SFT, not Pareto dominance.

\subsection{Required controls}

\begin{table}[H]
\centering
\caption{Minimum prospective comparison set.  ``Missing'' means no claim-bearing Paper B result
exists yet, not that an implementation can never be built.}
\label{tab:controls}
\small
\begin{tabularx}{\linewidth}{@{}lXl@{}}
\toprule
Arm & Purpose & Status \\
\midrule
Base & Single-pass pre-guard-SFT anchor & Retrospective scores ready \\
SFT & Specialized single-pass reference & Retrospective scores ready \\
Base+SFT output & Fixed primary candidate in \Cref{eq:composition} & Retrospective only \\
SFT+SFT output & Equal-inference-cost generic ensemble control; aggregate shared seed pairs at checkpoint level & Missing \\
WiSE-FT & Weight-space base/adapted interpolation control & Missing GPU rescoring \\
KL/replay tuning & Matched-compute preservation baseline & Missing training \\
\bottomrule
\end{tabularx}
\end{table}

Seed pairs in the SFT+SFT arm share adapters.  \PilotNSeedPairs{} unordered pairs from \PilotNSeeds{} seeds therefore
cannot be treated as \PilotNSeedPairs{} independent checkpoint replicates.  The locked analysis must aggregate
at checkpoint level and preserve the dependency induced by adapter reuse.  WiSE-FT and the
preservation-tuning arm need compatible retained weights or new training and rescoring.

\subsection{Heterogeneity and systems analysis}

The heterogeneity hypothesis requires a larger, lineage-diverse checkpoint panel.  Base
competence will be measured on a locked development cohort; the outcome will be prospective
composition-minus-SFT transfer macro-AP on a disjoint cohort.  Model identity and lineage must
remain visible rather than being hidden behind one correlation coefficient.

Every arm will report end-to-end latency, throughput, peak accelerator memory, training
accelerator-hours, and an inference-compute proxy on fixed hardware with a hashed measurement
protocol.  This separates matched training compute from matched inference compute.  Weight-space
composition may have a one-pass serving advantage; output composition and SFT+SFT are two-pass
systems and must not be described as cost-free.

\section{Implementation and reproducibility}
\label{sec:implementation}

The current workspace has two executable boundaries.  First,
\texttt{code/build\_pilot\_artifacts.py} validates the exact reviewed retrospective JSON and its
metadata, then generates all TeX macros and result tables plus a SHA-256 manifest.  Its check mode
fails when a generated file is missing or stale.  Second,
\texttt{code/validate\_protocol.py} validates the prospective contract.  Draft mode surfaces the
known missing fields; locked mode rejects them.

The current full cloud archive was checksum-verified during the Paper A audit but is not
distributed in this repository.  The public score cache can reverify the score and public-manifest
bindings, but it cannot independently rehash omitted adapter bytes and complete run metadata.
Third-party end-to-end re-verification therefore remains incomplete until those artifacts are
published at a durable immutable location.  This limitation is separate from, and does not excuse,
the need for an uninspected prospective Paper B cohort.

\section{Limitations and broader scope}
\label{sec:limitations}

The pilot includes \PilotNModels{} compact checkpoints from only two broad model lineages, one guard SFT
recipe, English input-prompt classification, benchmark-macro AP, and \PilotNAPDatasets{} AP datasets.  It does
not cover response moderation, tool calls, multi-turn dialogue, multilingual policies, arbitrary
new taxonomies, frontier hosted APIs, adversarial production traffic, or human escalation
workflows.  Benchmark label policies differ, and macro-averaging does not make them interchangeable.

The base is single-pass and lacks base-training-seed replication.  The bootstrap is conditional on
the selected panel.  Transfer rows were not used for SFT, but prior benchmark exposure prevents a
prospective interpretation.  The current operator was analyzed after Paper A's lock, and the
apparently strongest ablation was visible before this draft.  Threshold-free ranking improves
while the common-threshold diagnostic misses its calibration FPR target.  No standardized latency,
throughput, memory, or compute results exist yet.  These are study-design constraints, not caveats
that can be repaired with stronger prose.

\section{Conclusion}
\label{sec:conclusion}

In this retrospective fixed panel, a fixed average of separately calibrated base and SFT unsafe
scores exchanges lower represented-source AP for higher dataset-held-out transfer AP relative to
SFT.  The comparison with the base is heterogeneous, and the composition misses the transferred
FPR target.  The result is strong enough to justify a focused prospective study, but not strong
enough to establish dominance, a mechanism, a universal remedy, or a deployment recommendation.
The next scientific step is not a more confident headline.  It is a separately locked evaluation
on a genuinely uninspected cohort with an equal-cost SFT+SFT control, actual WiSE-FT rescoring,
matched-compute preservation tuning, independent competence measurement, and complete systems
costs.

\bibliographystyle{plainnat}
\bibliography{refs}

\appendix
\section{Retrospective publication contract}
\label{app:contract}

The manuscript-consumed evidence has status \PilotEvidenceStatus.  Its generated manifest records
the full composition, Paper A lock, score, and composition-analyzer hashes.  The TeX source contains
no manually entered pilot AP result.  The observed point contrasts are computed from the absolute
panel estimates, while bootstrap means are retained under separately named macros.  This prevents
the common reporting bug of presenting a bootstrap-resample mean as though it were the original
observed statistic.

The renderer is intentionally pinned to the reviewed retrospective object.  A changed analyzer or
new result must be reviewed and explicitly re-anchored; it cannot silently alter the draft.  The
future prospective pipeline should write into a separate Paper B artifact namespace rather than
overwriting this preview.

\section{Prospective completion checklist}
\label{app:checklist}

Before the draft can become a claim-bearing paper, the following items must all be complete:

\begin{enumerate}[leftmargin=1.5em]
  \item select, document, and hash a genuinely uninspected prospective cohort;
  \item justify the represented-source noninferiority margin without using prospective outcomes;
  \item freeze the operator hierarchy, model panel, estimands, multiplicity policy, RNGs, and
        family/seed resampling design;
  \item implement and verify SFT+SFT, WiSE-FT, and matched-compute KL/replay controls;
  \item measure base competence on a disjoint development cohort and expand the checkpoint panel;
  \item bind source, configuration, environment, container/runtime, and hardware measurement
        records in a separate Paper B lock;
  \item execute the locked analysis once, preserve failures, and report latency, throughput, peak
        memory, training compute, and inference compute for every arm; and
  \item publish sufficient immutable artifacts for independent verification under applicable data
        and model licenses.
\end{enumerate}

\end{document}
